\section{Kidz, Kulcha, Kreativity}

A few years ago, I watched a contrived controversy on a cable news channel. An exclusive and expensive liberal arts college in the suburbs of Los Angeles was giving students access to time on the campus cable TV channel. Not surprisingly, one group decided to produce pornography. Appearing on the show was a young woman, who was the producer, and a young man, one of the stars of the show.

Of course, the discussion followed the old tired pattern. On one side was the progressive liberated free thinkers and on the other the puritanical reactionary conservatives. Unfortunately, the world, even the educated world, can think only in terms of dualities when they need to think in terms of triads.

The format of the show was clearly designed to excite the two soul powers. One side represented liberated \emph{eros}, not subject to any constraint. The other side was incensed (\emph{thumos}) over the degeneration of youth. The missing term of the triad is, of course,
\emph{nous}, the intelligence.


The disappointment is that these privileged students could not conceive of anything higher or more creative than pornography, which doesn't take much intelligence at all. Suppose, instead, that had attempted some of these more creative projects.

\begin{itemize}
\item Weekly series based on stories from the Eddas

\item Plato's Republic has many story ideas

\begin{itemize}
\item A story based on the ideal republic (even has its prurient interest)
  
\item A story based on the allegory of the cave
\end{itemize}


\item Contests based on medieval knights. For example, climbing up the inside of a ladder in full armor.

\item Legends of Hyperborea, Agarttha, Rama, Orpheus, or the stories of other initiates
\end{itemize}


Perhaps projects like these have not been done because they can't be done. Perhaps something has been lost in the transition from an oral, to a written, to a visual society, so it is only the spectacle that captures attention, and not old myths with a metaphysical message.

\textbf{Leonardo da Vinci}, who knew something about the visual arts, preemptively rejected the portrayal of sexual activity.

\begin{quotex}
The art of procreation and the members employed therein are so repulsive, that if it were not for the beauty of the faces and the adornments of the actors and the pent-up impulse, nature would lose the human species.
\end{quotex}


Obviously, Leonardo did not consider the visual society, where pornography has become a major industry, and the repulsive is embraced.

In the piazzas of Florence where he grew up, he would have seen the puppet shows. In the evenings, the shows would have included sexual innuendos, in particular, mocking the affairs and trysts of the rich and famous. In the visual age, however, a staged photo or a leaked video on the Internet is often the path to stardom. With visual stimulation, the primal drives of eros (sex) and thumos (aggression) can be directly appealed to, as the content of television shows and movies attests to. There is an opposition that sees this, and its attendant dangers, but they have nothing to counter it. They appeal to sentimentalism, and say that man, on the contrary, is motivated by altruism, compassion, and so on. But there is not really a contrary, so sentimentalism is often disguised thumos, manifesting as the urge to appear as high status in another's eyes.

No, the resolution lies in \emph{nous}. However, \emph{nous} is non-formal manifestation; that is, there is no visual or other sensory representation of it. Hence, there is no direct way to portray it in any visual medium. But pornography is its opposite as Leonardo realized:

\begin{quotex}
Intellectual passion drives out sensuality. Whoso curbs not lustful desires puts himself on a level with the beasts.
\end{quotex}



\flrightit{Posted on 2011-01-26 by Cologero}

\begin{center}* * *\end{center}

\begin{footnotesize}
\begin{sffamily}
\texttt{James O'Meara on 2011-01-27 at 10:59 said:}

Gore Vidal, no prude he, said years ago that pornography was only for those without their own imaginations.

A recent book, Pornland, though somewhat tiresomely modernist-chauvinist, does make an interestingly similar point. `Pornography' today does not refer to say, D. H. Lawrence a Traditionalist at least in intention being banned as `indecent' but rather to the products of a billion or so dollar industry, which PROVIDES content, obviously with a certain agenda, to fill otherwise empty heads. One is not `free' but precisely brainwashed.

Now, these kids are likely a little more under their own power, since they are at least making their own stuff, but I would bet the actually content is the same. The projects you suggest would require an active imagination, which is precisely what is lacking today.

I notice that your suggestions seem to be based largely on oral traditions; Homer certainly, with Plato dialogues written down, perhaps to force the point a medieval epic read aloud to a courtly audience. Perhaps as Plato suggests, the move from oral to written involves a diminution of independence and creativity, and contrary to McLuhan, the move from print to visual yet more.

This would seem to be in accord with Traditional notions of esoteric ideas requiring oral transmission and, or because, requiring the student to assimilate them for himself. Evola talks of scientific knowledge being the sort that is `public' only because it is dead ideas poured into empty minds.

\hfill

\texttt{Perennial on 2011-01-27 at 12:32 said:}

Interesting thought, James. Makes one wonder... 

\hfill

\texttt{threefold\_crown on 2011-01-28 at 11:24 said:}

``The art of procreation and the members employed therein are so repulsive, that if it were not for the beauty of the faces and the adornments of the actors and the pent-up impulse, nature would lose the human species.''

Is this the L. da Vinci quote you refer to?

\hfill

\texttt{Cologero on 2011-01-28 at 12:06 said:}

Thank you, sir, for digging up that quote!

\hfill

\end{sffamily}
\end{footnotesize}

